\section{Progressive Conditioning for Diversity Measurement}

\subsection{Why Not Pairwise MI?}\label{sec:pairwise-mi}

A natural starting point is pairwise mutual information: measure how much seeing one response helps $\theta$ predict another.
This pairwise framing is also the one used by prior information-theoretic work on dialog diversity: Maximum Mutual Information (MMI) decoding \citep{li2016diversity} and Adversarial Information Maximization \citep{zhang2018aim} both optimize $I(X;Y)$ between input $X$ and a single response $Y$ at training or decoding time, to discourage generic replies.
Our construction differs from theirs in two further respects: we apply MI \emph{among responses from the same policy} rather than across the query–response boundary, and as a post-hoc evaluation diagnostic rather than a training objective.

Even setting those differences aside, pairwise MI among responses cannot detect a key failure mode of low diversity: when one response is a blend or interpolation of two others.
All three pairwise MI values may be near zero (the blend is not especially predictable from either parent alone), yet the set is clearly not diverse.
Progressive conditioning (Section~\ref{sec:progressive-conditioning}) handles this by conditioning on \emph{all} previous responses simultaneously, so the blend becomes predictable once both parents are in context.

\subsection{The Coherence Problem}\label{sec:coherence-problem}

Low MI (or low average PMI) means responses are approximately statistically independent under $\theta$'s assessment.
But there are two distinct ways to achieve independence:
\begin{enumerate}[nosep]
    \item \textbf{Diverse-coherent}: $\pi$ has broad support over plausible completions, each individually likely under $\theta$.
    \item \textbf{Random/incoherent}: $\pi$ produces noise; each sample is improbable under $\theta$ but trivially independent.
\end{enumerate}

Both yield low pairwise MI.
The distinguishing signal lies in the probability $\theta$ puts on the response without seeing other responses, $h_\theta(r \mid p)$, not in the pairwise structure.
This motivates decomposing the metric into separate coherence and independence components.\footnote{\citet{tevet2020evaluating} report that in pilot experiments, NLP graduate students rated sets of high-diversity but low-quality responses as low-diversity; they label this a human bias and design their crowdsourcing protocol to suppress it, but we read the same observation as evidence that coherence is part of what ``diversity'' of a generated set actually means.}

\subsection{Progressive Conditional Surprise}\label{sec:progressive-conditioning}

Rather than all-pairs PMI (a metric which cannot recognize the case when response 3 is a mix of response 1 and response 2 as low-diversity), we use the chain rule of mutual information.
Define:
\begin{equation}\label{eq:ak}
    a_k = -\log_2 \theta(r_k \mid r_{<k}, p) \quad \text{for } k = 1, \ldots, n
\end{equation}

Units: bits.
Each $a_k$ is the total surprise of the $k$-th response given all previous responses.
The sequence $(a_1, a_2, \ldots, a_n)$ is the \textbf{progressive conditional surprise curve}.
In practice, we normalize each $a_k$ by dividing by the byte count $\|r_k\|$ to obtain a per-byte curve (bits/byte); this makes the curve independent of $\theta$'s tokenizer and removes length variation.
Unless otherwise noted, all reported $a_k$ values and the diversity score $D_{Ca_n}$ are in per-byte units.

Since responses have different lengths, the raw total-bits $a_k$ values for a single ordering reflect both structural learning and length variation.
Per-byte normalization and averaging over random permutations of the response ordering (Section~\ref{sec:ordering}) remove these effects: each position averages over all responses, so the curve reflects only how $\theta$'s predictions improve with more context.

If $\theta$ were the true joint distribution of $\pi$'s outputs, then $\E[a_k]$ would be non-increasing in $k$ by the information-theoretic identity $H(X \mid Y, Z) \leq H(X \mid Y)$.
Since $a_k$ is a \emph{cross}-entropy under $\theta$ rather than the true conditional entropy, this monotonicity is an empirical property of $\theta$'s in-context learning, not a theorem.
In practice, we rely on $\theta$ being capable enough that the model-mismatch component does not dominate.
The monotonicity need not be perfect: mild violations (such as the last-position uptick discussed in Section~\ref{sec:tevet}, where \tevetLastUptickPct\% of Tevet samples have $a_5 > a_4$) add noise to the $a_n$ estimate but do not prevent the metric from discriminating diverse from non-diverse response sets.

The $a_k$ curve is the fundamental object of interest.
Its shape reveals the structure among the policy's outputs:
\begin{itemize}[nosep]
    \item For small $k$, each new response is likely to contain patterns $\theta$ has not yet seen, so $a_k$ remains near $a_1$.
In practice, cross-mode learning causes $a_k$ to drop even when responses are varied, but less than when they are similar.
    \item As $k$ grows, $\theta$ begins to recognize recurring structure (shared modes, stylistic regularities, topic patterns) and $a_k$ declines.
    \item For large $k$, $a_k$ converges to a floor $a_\infty > 0$ reflecting irreducible within-pattern variation that $\theta$ cannot predict from context alone.
\end{itemize}

This is analogous to a scree plot in PCA: the location of the ``elbow'' reveals when $\theta$ has learned most of the available structure, the floor corresponds to irreducible noise variance, and the initial height reflects the total variance.
Overlaying $a_k$ curves from different policies on the same axes provides the same kind of immediate visual diagnostic that comparing scree plots does.

\paragraph{Curve shape.} If modes were independent (seeing one mode's response tells $\theta$ nothing about other modes), the curve shape would depend on the ratio $m/n$: sigmoidal with a plateau at high $m/n$, exponential decay at low $m/n$.
Appendix~\ref{app:k0-derivation} derives this prediction formally.
In practice, capable base models exhibit cross-mode learning (Section~\ref{sec:cross-mode-learning}), which violates the independence assumption and produces exponential decay at all $m/n$ ratios.
The sigmoid prediction holds only for weaker models like GPT-2.

